
\section{Conclusion}

%In this paper, we introduced deep predictive learning which is the robot motion learning concept. We also summarized the background technology of predictive coding, related studies, and robot applications of our proposed concept. The three robot application examples described in terms of the composition of multiple motions are as follows; switching between multiple modules based on prediction errors, designing the dynamical systems inside the DNN model, and acquiring an integrated language--motion representation. In addition, we have shown their applicability to a wide variety of robotic tasks, including flexible object manipulation. Although each method has been developed separately, there are many common technical elements and it is expected that they will be integrated in the future. However, there are challenges to improving accuracy and applying the system to even more difficult tasks. We plan to revisit the proposed concept from the perspective of cognitive robotics, which is the origin of the concept, and apply the method to longer-term and more diverse tasks in line with real-world problems.
In this study, we introduce deep predictive learning, a concept of robot motion learning. We also summarize the background technology of predictive coding, related studies, and robotic applications of the proposed concept. The three robot application examples described regarding the composition of multiple motions are as follows: switching between multiple modules based on prediction errors, designing dynamic systems inside the DNN model, and acquiring an integrated language-motion representation. In addition, we demonstrate their applicability to various robotic tasks, including flexible object manipulation. Although each method has been developed separately, many common technical elements are expected to be integrated in the future. However, there are challenges in improving the accuracy and applying the system to more difficult tasks. We plan to revisit the proposed concept from the perspective of cognitive robotics the concept's origin, and apply the method to long-term and more diverse tasks in line with real-world problems.

